\chapter{Równania Sylvestera i Lapunowa}
\label{ch:lin22-sylvester}
\index{równanie Sylvestera@Równanie Sylvestera}
\index{równanie Lapunowa@Równanie Lapunowa}
\index{operator przeplatający@Operator przeplatający}
\index{stabilność spektralna@Stabilność spektralna}

\noindent\textbf{Cel rozdziału.}
W Rozdziale~\ref{ch:lin21-commutant}
opisaliśmy wszystkie operatory spełniające
równanie \(AT=TA\).
Zajmiemy się jego uogólnieniem:
równaniem jednorodnym \(AX=XB\)
i niejednorodnym \(AX-XB=C\).
Zobaczymy, że czynniki niezmiennicze
rozstrzygają istnienie i jednoznaczność
rozwiązania nad dowolnym ciałem.
Następnie, nad \(\mathbb C\), wykorzystamy
funkcje macierzy z
Rozdziału~\ref{ch:lin18-chevalley}
do rozwiązania równania Lapunowa
i opisu stabilności układu liniowego.

W pierwszych czterech sekcjach
\(K\) jest dowolnym ciałem;
\(A\in\operatorname{Mat}_{n\times n}(K)\)
i \(B\in\operatorname{Mat}_{m\times m}(K)\),
gdzie \(n,m\ge1\).
Macierz \(X\) ma rozmiar \(n\times m\).
W części o równaniu Lapunowa
przyjmujemy jawnie \(K=\mathbb C\).

\section{Przeplatanie operatorów}
\label{sec:lin22-intertwine}

\begin{definicja}[Równanie Sylvestera]
\label{def:lin22-sylvester}
Dla ustalonych macierzy
\(A\in\operatorname{Mat}_{n\times n}(K)\)
i \(B\in\operatorname{Mat}_{m\times m}(K)\)
definiujemy liniowy operator Sylvestera
\[
 \mathcal S_{A,B}\colon
 \operatorname{Mat}_{n\times m}(K)
 \longrightarrow\operatorname{Mat}_{n\times m}(K),
 \qquad \mathcal S_{A,B}(X)=AX-XB.
\]
Równanie
\(\mathcal S_{A,B}(X)=C\),
gdzie \(C\in\operatorname{Mat}_{n\times m}(K)\),
nazywamy \emph{równaniem Sylvestera}.
Równanie \(AX=XB\) nazywamy
\emph{równaniem przeplatania},
a jego rozwiązania
\emph{operatorami przeplatającymi}.
\end{definicja}

\begin{stwierdzenie}[Przestrzeń operatorów przeplatających]
\label{prop:lin22-hom}
Niech \(V_A\) i \(W_B\) będą
modułami nad \(K[t]\), otrzymanymi
z działania macierzy \(A\) na \(K^n\)
i macierzy \(B\) na \(K^m\),
jak w Definicji~\ref{def:lin19-module}.
Wówczas
\[
 \ker\mathcal S_{A,B}
 =\operatorname{Hom}_{K[t]}(W_B,V_A),
\]
przy utożsamieniu macierzy \(X\)
z odwzorowaniem \(K\)-liniowym
\(K^m\to K^n\).
Jeżeli \(X_0\) rozwiązuje
\(AX-XB=C\), to zbiór wszystkich
jego rozwiązań jest równy
\(X_0+\ker\mathcal S_{A,B}\).
\end{stwierdzenie}
\begin{proof}
Macierz \(X\) należy do jądra
\(\mathcal S_{A,B}\) dokładnie wtedy,
gdy \(AX=XB\).
Indukcja po \(k\ge0\) daje wtedy
\(A^kX=XB^k\), a liniowe sumowanie
daje \(p(A)X=Xp(B)\)
dla każdego \(p\in K[t]\).
Jest to dokładnie warunek
\(K[t]\)-liniowości odwzorowania
z \(W_B\) do \(V_A\).
Odwrotnie, zachowanie działania
wielomianu \(t\) daje \(AX=XB\).

Jeżeli \(\mathcal S_{A,B}(X_0)=C\),
to \(\mathcal S_{A,B}(X)=C\)
jest równoważne
\(\mathcal S_{A,B}(X-X_0)=0\),
ponieważ \(\mathcal S_{A,B}\)
jest liniowy.
Zatem wszystkie rozwiązania
mają postać \(X_0+Y\),
\(Y\in\ker\mathcal S_{A,B}\).
\end{proof}

\section{Kryterium jednoznacznej rozwiązywalności}
\label{sec:lin22-criterion}

\begin{twierdzenie}[Wymiar przestrzeni przeplatań]
\label{thm:lin22-kernel}
Niech \(f_1\mid\cdots\mid f_r\)
oraz \(g_1\mid\cdots\mid g_s\)
oznaczają odpowiednio moniczne
czynniki niezmiennicze macierzy
\(A\in\operatorname{Mat}_{n\times n}(K)\)
i \(B\in\operatorname{Mat}_{m\times m}(K)\).
Wówczas
\[
 \dim_K\ker\mathcal S_{A,B}
 =\sum_{i=1}^r\sum_{j=1}^s
          \deg\gcd(f_i,g_j).
\]
W szczególności przestrzeń
operatorów przeplatających jest
zerowa dokładnie wtedy, gdy
wielomiany charakterystyczne
\(\chi_A\) i \(\chi_B\)
są względnie pierwsze.
\end{twierdzenie}
\begin{proof}
Z Twierdzenia~\ref{thm:lin19-frobenius}
otrzymujemy izomorfizmy modułów
\[
 V_A\cong\bigoplus_{i=1}^r K[t]/(f_i),
 \qquad
 W_B\cong\bigoplus_{j=1}^s K[t]/(g_j).
\]
Każdy homomorfizm z drugiej
sumy do pierwszej jest jednoznacznie
wyznaczony przez składowe
\(\varphi_{ij}\colon K[t]/(g_j)
\to K[t]/(f_i)\), uzyskane przez
włożenia i rzuty na składniki.
Odwrotnie, dowolna rodzina takich
homomorfizmów określa odwzorowanie
sum przez zsumowanie składowych.
Zatem przestrzeń Hom jest
sumą prostą przestrzeni tych składowych.
Na mocy
Twierdzenia~\ref{thm:lin21-hom-cyclic}
jej wymiar jest sumą
\(\deg\gcd(f_i,g_j)\).
To daje pierwszy wzór dzięki
Stwierdzeniu~\ref{prop:lin22-hom}.

Wszystkie składniki tej sumy
są nieujemne.
Wymiar jądra jest zerowy
dokładnie wtedy, gdy każda para
\(f_i,g_j\) jest względnie pierwsza.
Ponieważ
\(\chi_A=\prod_i f_i\) i
\(\chi_B=\prod_jg_j\)
na mocy
Twierdzenia~\ref{thm:lin19-frobenius},
warunek ten jest równoważny
\(\gcd(\chi_A,\chi_B)=1\).
Istotnie, wspólny nierozkładalny
dzielnik dwóch iloczynów
dzieli co najmniej jeden czynnik
każdego z nich, na mocy
lematu Euklidesa w \(K[t]\).
\end{proof}

\begin{twierdzenie}[Kryterium Sylvestera]
\label{thm:lin22-sylvester}
Niech \(K\) będzie dowolnym ciałem,
\(n,m\ge1\), a \(A\) i \(B\)
macierzami kwadratowymi rozmiarów
odpowiednio \(n\) i \(m\)
o współczynnikach w \(K\).
Następujące warunki są równoważne:
\begin{enumerate}[label=\textup{(\roman*)}]
 \item dla każdego
 \(C\in\operatorname{Mat}_{n\times m}(K)\)
 równanie \(AX-XB=C\)
 ma dokładnie jedno rozwiązanie \(X\);
 \item operator \(\mathcal S_{A,B}\)
 jest odwracalny;
 \item jedynym rozwiązaniem
 równania \(AX=XB\) jest \(X=0\);
 \item \(\gcd(\chi_A,\chi_B)=1\).
\end{enumerate}
Jeżeli \(K=\mathbb C\), warunek
\textup{(iv)} jest równoważny
rozłączności zbiorów wartości
własnych \(A\) i \(B\).
\end{twierdzenie}
\begin{proof}
Z Definicji~\ref{def:lin22-sylvester}
operator \(\mathcal S_{A,B}\)
jest endomorfizmem
\(nm\)-wymiarowej przestrzeni
\(\operatorname{Mat}_{n\times m}(K)\).
Warunek (i) mówi dokładnie,
że jest on bijektywny, czyli (ii).
Dla endomorfizmu skończenie wymiarowej
przestrzeni warunki odwracalności
i zerowości jądra są równoważne:
z twierdzenia o wymiarze
obrazu i jądra wynika,
że injektywność wymusza
surjektywność i odwrotnie.
To daje równoważność (ii) i (iii).

Z Twierdzenia~\ref{thm:lin22-kernel}
jądro operatora Sylvestera
jest zerowe dokładnie dla
\(\gcd(\chi_A,\chi_B)=1\),
czyli (iv).
Nad \(\mathbb C\) oba wielomiany
charakterystyczne rozkładają się
na czynniki liniowe, a ich
moniczne wspólne nierozkładalne
czynniki odpowiadają dokładnie
wspólnym wartościom własnym.
To daje ostatnią równoważność.
\end{proof}

\begin{wniosek}[Wymiar rodziny rozwiązań]
\label{cor:lin22-affine}
Niech \(A,B\) będą macierzami
jak w Twierdzeniu~\ref{thm:lin22-sylvester}
i niech
\(C\in\operatorname{Mat}_{n\times m}(K)\).
Jeśli równanie \(AX-XB=C\)
ma choć jedno rozwiązanie,
to zbiór wszystkich jego rozwiązań
jest podprzestrzenią afiniczną
o wymiarze
\[
 \sum_{i=1}^r\sum_{j=1}^s
       \deg\gcd(f_i,g_j),
\]
gdzie \(f_i,g_j\) są czynnikami
niezmienniczymi \(A,B\).
Jeżeli ten wymiar jest równy zeru,
równanie ma dokładnie jedno rozwiązanie
dla każdego \(C\).
\end{wniosek}
\begin{proof}
Z
Stwierdzenia~\ref{prop:lin22-hom}
zbiór rozwiązań ma postać
\(X_0+\ker\mathcal S_{A,B}\),
a przesunięcie nie zmienia wymiaru.
Jego wymiar wyznacza
Twierdzenie~\ref{thm:lin22-kernel}.
Jeśli jądro jest zerowe,
z Twierdzenia~\ref{thm:lin22-sylvester}
wynika również istnienie
rozwiązania dla każdego \(C\).
\end{proof}

\begin{przyklad}[To samo spektrum: brak lub nadmiar rozwiązań]
\label{prz:lin22-obstruction}
Niech \(A=B=(2)\), gdzie \(K=\mathbb Q\).
Wtedy \(\mathcal S_{A,B}(x)=2x-2x=0\)
dla każdego \(x\in\mathbb Q\).
Równanie \(2x-2x=1\)
nie ma rozwiązania,
natomiast równanie
\(2x-2x=0\) ma nieskończenie
wiele rozwiązań.
Wielomiany charakterystyczne
obu macierzy są równe \(t-2\),
więc nie są względnie pierwsze.
Wzór z
Twierdzenia~\ref{thm:lin22-kernel}
daje wymiar jądra równy \(1\),
zgodnie z obliczeniem.
\end{przyklad}

\section{Wzór wykonawczy na rozwiązanie}
\label{sec:lin22-explicit}

\begin{lemat}[Tożsamość teleskopowa dla Sylvestera]
\label{lem:lin22-telescope}
Dla dowolnych \(A\in\operatorname{Mat}_{n\times n}(K)\),
\(B\in\operatorname{Mat}_{m\times m}(K)\),
\(X\in\operatorname{Mat}_{n\times m}(K)\)
oraz liczby całkowitej \(k\ge1\)
zachodzi
\[
 A^kX-XB^k
 =\sum_{j=0}^{k-1}
 A^{k-1-j}(AX-XB)B^j.
\]
\end{lemat}
\begin{proof}
Po wymnożeniu prawej strony
otrzymujemy
\[
 \sum_{j=0}^{k-1}A^{k-j}XB^j
 -\sum_{j=0}^{k-1}A^{k-1-j}XB^{j+1}.
\]
Po zmianie indeksu \(l=j+1\)
w drugiej sumie wszystkie składniki
dla indeksów od \(1\) do \(k-1\)
skracają się parami.
Pozostają \(A^kX\) z pierwszej
sumy i \(-XB^k\) z drugiej.
\end{proof}

\begin{twierdzenie}[Rozwiązanie przez wielomian charakterystyczny]
\label{thm:lin22-explicit}
Niech \(A\in\operatorname{Mat}_{n\times n}(K)\),
\(B\in\operatorname{Mat}_{m\times m}(K)\)
z \(n,m\ge1\), przy czym
\(\gcd(\chi_A,\chi_B)=1\).
Zapiszmy
\[
 h(t)=\chi_A(t)=\sum_{k=0}^{n}h_kt^k.
\]
Dla każdego
\(C\in\operatorname{Mat}_{n\times m}(K)\)
jedynym rozwiązaniem
\(AX-XB=C\) jest
\[
 X=-\left(
 \sum_{k=1}^{n}h_k
 \sum_{j=0}^{k-1}
   A^{k-1-j}CB^j
 \right)h(B)^{-1}.
\]
Macierz \(h(B)\) jest odwracalna
na mocy założeń.
\end{twierdzenie}
\begin{proof}
Ponieważ \(\gcd(h,\chi_B)=1\),
tożsamość Bézouta daje
\(u(t)h(t)+v(t)\chi_B(t)=1\)
dla pewnych \(u,v\in K[t]\).
Podstawiając \(B\) oraz korzystając
z twierdzenia Cayleya--Hamiltona
z Rozdziału~\ref{ch:lin16-niezmiennicze},
otrzymujemy
\(u(B)h(B)=I_m\).
Zatem \(h(B)\) ma lewą odwrotność,
a jako macierz kwadratowa nad ciałem
jest odwracalna.

Niech \(X\) będzie dowolnym
rozwiązaniem rozpatrywanego równania.
Lemat~\ref{lem:lin22-telescope}
daje dla \(k\ge1\)
\[
 A^kX-XB^k
 =\sum_{j=0}^{k-1}A^{k-1-j}CB^j.
\]
Mnożąc przez \(h_k\)
i sumując dla \(k=1,\ldots,n\),
możemy dodać człon \(k=0\),
który jest zerowy.
Stąd
\[
 h(A)X-Xh(B)=
 \sum_{k=1}^{n}h_k
 \sum_{j=0}^{k-1}A^{k-1-j}CB^j.
\]
Twierdzenie Cayleya--Hamiltona
daje \(h(A)=0\),
a pomnożenie ostatniej równości
z prawej strony przez \(h(B)^{-1}\)
prowadzi do podanego wzoru.

Istnienie i jednoznaczność
rozwiązania dla każdego \(C\)
wynikają niezależnie z
Twierdzenia~\ref{thm:lin22-sylvester}.
Ponieważ wzór jest konieczny
dla tego jedynego rozwiązania,
rzeczywiście je wyznacza.
\end{proof}

\begin{przyklad}[Równanie z macierzą trójkątną]
\label{prz:lin22-explicit}
Niech \(K=\mathbb Q\),
\[
 A=\begin{pmatrix}1&1\\0&1\end{pmatrix},
 \quad B=(3),
 \quad C=\begin{pmatrix}2\\4\end{pmatrix}.
\]
Równanie \(AX-XB=C\)
ma postać
\[
 \begin{pmatrix}-2&1\\0&-2\end{pmatrix}
 \begin{pmatrix}x\\y\end{pmatrix}
 =\begin{pmatrix}2\\4\end{pmatrix}.
\]
Drugie równanie daje \(y=-2\),
a pierwsze \(-2x+y=2\),
więc \(x=-2\).
Rozwiązaniem jest zatem
\(X=(-2,-2)^{\mathsf T}\).

Sprawdzimy to także z
Twierdzenia~\ref{thm:lin22-explicit}.
Ponieważ \(h(t)=(t-1)^2=t^2-2t+1\),
mamy
\[
 \sum_{k=1}^{2}h_k
 \sum_{j=0}^{k-1}A^{k-1-j}CB^j
 =-2C+AC+CB.
\]
Tutaj \(AC=(6,4)^{\mathsf T}\),
\(CB=(6,12)^{\mathsf T}\),
więc suma wynosi
\((8,8)^{\mathsf T}\).
Z kolei \(h(B)=(3-1)^2=4\),
stąd \(X=-(8,8)^{\mathsf T}/4\),
jak poprzednio.
\end{przyklad}

\section{Dwa bloki Jordana}
\label{sec:lin22-jordan}

\begin{stwierdzenie}[Wzór skończony dla różnych wartości własnych]
\label{prop:lin22-jordan-solve}
Niech \(\lambda,\mu\in K\),
\(\lambda\ne\mu\), oraz
\(A=J_a(\lambda)\),
\(B=J_b(\mu)\), gdzie \(a,b\ge1\).
Połóżmy
\(N=J_a(0)\), \(M=J_b(0)\)
oraz
\[
 \mathcal N(X)=NX-XM
 \quad\text{dla }X\in\operatorname{Mat}_{a\times b}(K).
\]
Wówczas \(\mathcal N^{a+b-1}=0\),
a jedyne rozwiązanie
\(AX-XB=C\) jest równe
\[
 X=\sum_{k=0}^{a+b-2}
 \frac{(-1)^k}{(\lambda-\mu)^{k+1}}\,
 \mathcal N^k(C).
\]
Wzór obowiązuje również
w dodatniej charakterystyce ciała.
\end{stwierdzenie}
\begin{proof}
Z definicji bloków Jordana
wynika
\[
 AX-XB=(\lambda-\mu)X+\mathcal N(X).
\]
Niech \(L(X)=NX\) i \(R(X)=XM\).
Wówczas \(L\) i \(R\) są
przemiennymi endomorfizmami przestrzeni
macierzy \(a\times b\), ponieważ
\(L(R(X))=NXM=R(L(X))\).
Ponadto \(L^a=0\) i \(R^b=0\).
Dla \(k=a+b-1\) wzór dwumianowy
na przemienne operatory daje
\[
 (L-R)^k
 =\sum_{j=0}^{k}{k\choose j}
       (-1)^jL^{k-j}R^j.
\]
W każdym składniku zachodzi
\(k-j\ge a\) albo \(j\ge b\):
gdyby oba warunki były fałszywe,
mielibyśmy
\(k-j\le a-1\) i \(j\le b-1\),
a więc \(k\le a+b-2\),
wbrew wybranej wartości \(k\).
Wszystkie składniki znikają,
stąd \(\mathcal N^k=0\).
Dwumianowy rachunek używa jedynie
całkowitych współczynników,
więc jest prawdziwy nad każdym ciałem.

Oznaczmy \(\delta=\lambda-\mu\ne0\).
Ponieważ \(\mathcal N^{a+b-1}=0\),
skończony szereg geometryczny
\[
 (\delta I+\mathcal N)^{-1}
 =\delta^{-1}\sum_{k=0}^{a+b-2}
           (-\delta^{-1}\mathcal N)^k
\]
jest odwrotnością operatora
\(\delta I+\mathcal N\):
przy mnożeniu kolejne składniki
skracają się, a pozostała
najwyższa potęga \(\mathcal N\)
jest zerowa.
Zastosowanie odwrotności do \(C\)
daje podany wzór.
\end{proof}

\begin{stwierdzenie}[Wspólna wartość własna]
\label{prop:lin22-jordan-hom}
Niech \(a,b\ge1\) oraz
\(\lambda,\mu\in K\).
Wtedy
\[
 \dim_K\{X\in\operatorname{Mat}_{a\times b}(K):
          J_a(\lambda)X=XJ_b(\mu)\}
 =\begin{cases}
  \min(a,b),&\lambda=\mu,\\
  0,&\lambda\ne\mu.
 \end{cases}
\]
\end{stwierdzenie}
\begin{proof}
Na mocy
Stwierdzenia~\ref{prop:lin20-linear-primary}
blok Jordana \(J_a(\lambda)\)
jest podobny do macierzy
towarzyszącej
\(C((t-\lambda)^a)\);
analogicznie dla \(J_b(\mu)\).
Odpowiednie moduły są zatem
izomorficzne z
\(K[t]/((t-\lambda)^a)\)
i \(K[t]/((t-\mu)^b)\).
Z Twierdzenia~\ref{thm:lin21-hom-cyclic}
wymiar przestrzeni przeplatań
jest stopniem największego
wspólnego dzielnika tych wielomianów.
Dla różnych \(\lambda,\mu\)
NWD wynosi \(1\);
dla \(\lambda=\mu\)
wynosi \((t-\lambda)^{\min(a,b)}\).
To daje obie wartości.
\end{proof}

\section{Równanie Lapunowa i stabilność}
\label{sec:lin22-lyapunov}

\begin{definicja}[Stabilność spektralna]
\label{def:lin22-hurwitz}
Niech \(A\in\operatorname{Mat}_{n\times n}(\mathbb C)\).
Macierz \(A\) nazywamy
\emph{spektralnie stabilną w czasie ciągłym}
(macierzą Hurwitza), jeżeli
dla każdej jej wartości własnej
\(\lambda\in\mathbb C\) zachodzi
\(\operatorname{Re}\lambda<0\).
\end{definicja}

\begin{definicja}[Równanie Lapunowa]
\label{def:lin22-lyapunov}
Dla macierzy
\(A,Q\in\operatorname{Mat}_{n\times n}(\mathbb C)\)
\emph{ciągłym równaniem Lapunowa}
nazywamy
\[
 A^*P+PA=-Q
\]
o niewiadomej
\(P\in\operatorname{Mat}_{n\times n}(\mathbb C)\),
gdzie \(A^*\) oznacza macierz
sprzężoną i transponowaną.
Macierz hermitowską \(H\)
nazywamy \emph{dodatnio określoną},
jeśli \(x^*Hx>0\)
dla każdego niezerowego
\(x\in\mathbb C^n\).
\end{definicja}

\begin{twierdzenie}[Twierdzenie Lapunowa o stabilności]
\label{thm:lin22-lyapunov}
Niech \(n\ge1\) i
\(A\in\operatorname{Mat}_{n\times n}(\mathbb C)\).
Następujące warunki są równoważne:
\begin{enumerate}[label=\textup{(\roman*)}]
 \item macierz \(A\) jest
 spektralnie stabilna w czasie ciągłym;
 \item dla każdej hermitowskiej,
 dodatnio określonej macierzy
 \(Q\in\operatorname{Mat}_{n\times n}(\mathbb C)\)
 równanie \(A^*P+PA=-Q\)
 ma dokładnie jedno rozwiązanie,
 a rozwiązanie to jest hermitowskie
 i dodatnio określone;
 \item istnieją hermitowskie,
 dodatnio określone macierze \(P,Q\)
 spełniające
 \(A^*P+PA=-Q\).
\end{enumerate}
Przy spełnionych warunkach
\textup{(i)--(iii)} rozwiązanie
dla każdego takiego \(Q\)
jest dane całką zbieżną bezwzględnie
\[
 P=\int_0^\infty e^{tA^*}Qe^{tA}\,\mathrm dt.
\]
\end{twierdzenie}
\begin{proof}
\emph{Implikacja (i) \(\Rightarrow\) (ii).}
Niech \(A\) będzie spektralnie stabilna,
a \(Q\) hermitowska i dodatnio określona.
Wykażemy najpierw zbieżność całki.
Nad \(\mathbb C\) istnieje
macierz odwracalna \(S\)
i suma bloków Jordana \(J\),
dla których \(A=SJS^{-1}\)
na mocy
Twierdzenia o postaci Jordana
z Rozdziału~\ref{ch:lin17-jordan}.
Na bloku \(J_d(\lambda)=\lambda I+N\),
gdzie \(N^d=0\),
rachunek wykładniczy
z Rozdziału~\ref{ch:lin18-chevalley}
daje
\[
 e^{tJ_d(\lambda)}
 =e^{t\lambda}
  \sum_{j=0}^{d-1}\frac{t^j}{j!}N^j
 \qquad(t\ge0).
\]
Ponieważ zbiór wartości własnych
jest skończony, istnieje \(\beta>0\)
takie, że
\(\operatorname{Re}\lambda\le-\beta\)
dla każdej wartości własnej.
W ustalonej podmultiplikatywnej
normie macierzy istnieje zatem
stała \(M>0\) taka, że
\[
 \|e^{tA}\|
 \le M(1+t^{n-1})e^{-\beta t}
 \qquad(t\ge0).
\]
Rzeczywiście, wystarczy oszacować
skończone sumy potęg \(t^j\)
na każdym bloku, a następnie
zastosować podobieństwo przez \(S\).
Ponadto
\(\|e^{tA^*}\|\) podlega
analogicznemu oszacowaniu,
gdyż \(e^{tA^*}=(e^{tA})^*\).
Dla normy operatorowej pochodzącej
od standardowego iloczynu
hermitowskiego mamy nawet równość
tych norm.
Wobec tego
\(\|e^{tA^*}Qe^{tA}\|\)
jest ograniczona przez stałą
wielokrotność
\((1+t^{n-1})^2e^{-2\beta t}\).
Ta funkcja jest całkowalna
na \([0,\infty)\), ponieważ
każdy jej składnik jest
wielomianem razy malejąca
funkcja wykładnicza.
Całka macierzowa jest więc
zbieżna bezwzględnie
w dowolnej normie
na skończenie wymiarowej przestrzeni
macierzy; można całkować
każdy jej współczynnik oddzielnie.

Połóżmy
\(P=\int_0^\infty e^{tA^*}Qe^{tA}\,\mathrm dt\).
Ponieważ \(Q^*=Q\),
każda macierz podcałkowa
jest hermitowska,
więc \(P^*=P\).
Dla niezerowego \(x\in\mathbb C^n\)
mamy
\[
 x^*Px=\int_0^\infty
       (e^{tA}x)^*Q(e^{tA}x)\,\mathrm dt.
\]
Funkcja podcałkowa jest ciągła
i przy \(t=0\) ma wartość
\(x^*Qx>0\).
Pozostaje więc dodatnia
na pewnym przedziale
\([0,\varepsilon]\), skąd
całka jest ściśle dodatnia.
Zatem \(P\) jest dodatnio określona.

Rozważmy
\(F(t)=e^{tA^*}Qe^{tA}\).
Różniczkując iloczyn i korzystając
z przemienności \(A\)
ze swą wykładniczą,
otrzymujemy
\[
 F'(t)=A^*F(t)+F(t)A.
\]
Wszystkie składniki są całkowalne;
możemy zatem całkować tę równość
wyraz po wyrazie na \([0,R]\):
\[
 A^*\int_0^R F(t)\,\mathrm dt
 +\left(\int_0^R F(t)\,\mathrm dt\right)A
 =F(R)-F(0).
\]
Z oszacowania wykładniczego
wynika \(F(R)\to0\), gdy \(R\to\infty\),
zaś \(F(0)=Q\).
Przechodząc do granicy,
dostajemy \(A^*P+PA=-Q\).

Aby wykazać jednoznaczność,
niech \(P_1,P_2\) będą dowolnymi
rozwiązaniami, bez założenia,
że są hermitowskie.
Różnica \(H=P_1-P_2\)
spełnia \(A^*H+HA=0\).
Wtedy dla
\(G(t)=e^{tA^*}He^{tA}\)
otrzymujemy
\(G'(t)=e^{tA^*}(A^*H+HA)e^{tA}=0\).
Zatem \(G(t)=G(0)=H\)
dla każdego \(t\ge0\).
Ponieważ \(e^{tA}\to0\)
i \(e^{tA^*}\to0\),
mamy jednocześnie \(G(t)\to0\).
Wynika stąd \(H=0\),
czyli rozwiązanie jest jedyne.
Dowodzi to (ii).

Implikacja (ii) \(\Rightarrow\) (iii)
wynika z wyboru \(Q=I_n\).

\emph{Implikacja (iii) \(\Rightarrow\) (i).}
Załóżmy, że hermitowskie
dodatnio określone \(P,Q\)
spełniają \(A^*P+PA=-Q\).
Niech \(\lambda\) będzie dowolną
wartością własną \(A\),
a \(v\ne0\) wektorem takim,
że \(Av=\lambda v\).
Po pomnożeniu równania Lapunowa
z lewej przez \(v^*\),
a z prawej przez \(v\),
otrzymujemy
\[
 -v^*Qv
 =v^*A^*Pv+v^*PAv
 =(\overline\lambda+\lambda)v^*Pv
 =2\operatorname{Re}\lambda\,v^*Pv.
\]
Obie liczby \(v^*Pv\)
i \(v^*Qv\) są dodatnie.
Stąd
\(\operatorname{Re}\lambda
 =-v^*Qv/(2v^*Pv)<0\).
Ponieważ \(\lambda\) była dowolna,
macierz \(A\) jest spektralnie stabilna.
\end{proof}

\begin{wniosek}[Malejąca funkcja kwadratowa]
\label{cor:lin22-energy}
Niech \(A\in\operatorname{Mat}_{n\times n}(\mathbb C)\)
będzie spektralnie stabilna,
a \(P,Q\) niech będą dodatnio
określonymi macierzami z
Twierdzenia~\ref{thm:lin22-lyapunov}.
Dla rozwiązania
\(x(t)=e^{tA}x_0\)
układu równań różniczkowych
\(x'(t)=Ax(t)\)
funkcja rzeczywista
\(E(t)=x(t)^*Px(t)\)
spełnia
\[
 E'(t)=-x(t)^*Qx(t).
\]
Jeżeli \(x_0\ne0\),
to \(E'(t)<0\)
dla każdego \(t\ge0\).
\end{wniosek}
\begin{proof}
Z pochodnej iloczynu
i równości \(x'=Ax\)
otrzymujemy
\[
 E'(t)
 =(Ax(t))^*Px(t)+x(t)^*PAx(t)
 =x(t)^*(A^*P+PA)x(t)
 =-x(t)^*Qx(t).
\]
Macierz \(e^{tA}\) jest odwracalna,
bo jej odwrotnością jest \(e^{-tA}\),
co wynika z własności wykładniczej
w Rozdziale~\ref{ch:lin18-chevalley}.
Zatem \(x_0\ne0\)
implikuje \(x(t)\ne0\).
Dodatnia określoność \(Q\)
daje wtedy \(E'(t)<0\).
\end{proof}

\begin{przyklad}[Równanie skalarne Lapunowa]
\label{prz:lin22-scalar}
Dla \(A=(-a)\), gdzie \(a>0\)
jest liczbą rzeczywistą,
i \(Q=(q)\), gdzie \(q>0\),
równanie Lapunowa ma postać
\[
 (-a)p+p(-a)=-q,
 \qquad p=\frac{q}{2a}>0.
\]
Wzór całkowy daje ten sam wynik:
\[
 \int_0^\infty e^{-at}q e^{-at}\,\mathrm dt
 =q\int_0^\infty e^{-2at}\,\mathrm dt
 =\frac q{2a}.
\]
\end{przyklad}

\begin{przyklad}[Równanie Lapunowa dla macierzy niediagonalnej]
\label{prz:lin22-non-normal}
Niech
\[
 A=\begin{pmatrix}-1&1\\0&-2\end{pmatrix},
 \qquad Q=I_2.
\]
Macierz \(A\) nie jest normalna,
lecz jej wartości własne
to \(-1\) i \(-2\),
więc jest spektralnie stabilna.
Szukamy rzeczywistej
macierzy symetrycznej
\(P=\begin{pmatrix}u&v\\v&w\end{pmatrix}\).
Bezpośrednie mnożenie daje
\[
 A^{\mathsf T}P+PA
 =\begin{pmatrix}
   -2u&u-3v\\
   u-3v&2v-4w
  \end{pmatrix}.
\]
Porównanie z \(-I_2\)
prowadzi do trzech równań
\(-2u=-1\), \(u-3v=0\),
\(2v-4w=-1\).
Stąd
\[
 P=\begin{pmatrix}
 \frac12&\frac16\\
 \frac16&\frac13
 \end{pmatrix}.
\]
Macierz jest dodatnio określona:
jej pierwszy główny minor wynosi
\(1/2>0\), a wyznacznik
\(1/6-1/36=5/36>0\),
więc z kryterium Sylvestera
dla form hermitowskich
wynika dodatnia określoność.
Można również zapisać wprost
\[
 (x,y)P(x,y)^{\mathsf T}
 =\frac12\left(x+\frac y3\right)^2
  +\frac5{18}y^2>0
\]
dla niezerowej rzeczywistej pary
\((x,y)\); w przypadku zespolonym
należy zastąpić kwadraty
modułami kwadratów i odpowiednio
zastosować sprzężenie.
\end{przyklad}

\begin{przyklad}[Brak dodatniego rozwiązania przy niestabilności]
\label{prz:lin22-unstable}
Niech
\(A=\operatorname{diag}(-1,1)\)
oraz \(Q=I_2\).
Wśród diagonalnych macierzy
\(P=\operatorname{diag}(u,w)\)
równanie Lapunowa wymaga
\(-2u=-1\) i \(2w=-1\),
czyli \(u=1/2\) i \(w=-1/2\).
Zatem takie rozwiązanie
nie jest dodatnio określone.
W rzeczywistości nie istnieje
\emph{żadne} dodatnio określone
rozwiązanie: dla wektora własnego
\(e_2\) odpowiadającego wartości
własnej \(1\) argument z dowodu
Twierdzenia~\ref{thm:lin22-lyapunov}
wymuszałby
\(2e_2^*Pe_2=-e_2^*Qe_2=-1\),
co przeczy dodatniej określoności \(P\).
Równanie nie ma również rozwiązania
jednoznacznego, ponieważ
w macierzy \(P=(p_{ij})\)
elementy pozadiagonalne
są mnożone przez sumę
wartości własnych \(-1+1=0\).
\end{przyklad}

\section{Podsumowanie}
\label{sec:lin22-summary}

Przestrzeń rozwiązań równania
\(AX=XB\) jest przestrzenią
homomorfizmów modułów
nad \(K[t]\). Jej wymiar
obliczamy z największych
wspólnych dzielników
czynników niezmienniczych
macierzy \(A,B\).
Równanie Sylvestera
\(AX-XB=C\)
ma dokładnie jedno rozwiązanie
dla każdego \(C\) wtedy
i tylko wtedy, gdy
\(\gcd(\chi_A,\chi_B)=1\);
w takim przypadku podaliśmy
skończony wzór wykonawczy.
Nad \(\mathbb C\) oznacza to
brak wspólnych wartości własnych.

Równanie Lapunowa
\(A^*P+PA=-Q\)
wiąże to zagadnienie
z dynamiką \(x'=Ax\).
Spektralna stabilność macierzy
jest równoważna istnieniu
dodatnio określonego rozwiązania
dla pewnej dodatnio określonej
prawej strony, a wtedy dla
każdej takiej prawej strony
istnieje dokładnie jedno rozwiązanie,
dane całką z wykładniczych macierzy.
